\section{Method}

We conducted a two-wave longitudinal survey study to examine the role of perceived economic inequality in predicting changes in students’ identification with educational values (IEV) and academic prosocial behavior (APB). Standardized questionnaires were administered in classroom settings, and data were analyzed using latent change score models.

\subsection{Participants}
Participants were 634 students (approximately 10 years old) recruited from a primary school in northern China. Data were collected in two waves, separated by a seven-month interval. Only students with valid responses at both waves were included in the final sample.

\subsection{Materials and Procedure}
Students completed paper-and-pencil questionnaires during regular school hours. Measures included perceived economic inequality, identification with educational values (IEV), academic prosocial behavior (APB), and subjective well-being. Data were analyzed using latent change score (LCS) models in Mplus, with measurement invariance tested across waves. Model fit was evaluated using multiple indices, including RMSEA, CFI, and SRMR.